\section{Discussion}

\plat{} suggests that culturally localized programming-language design is not
only a question of translating keywords. It requires decisions about who the
language is for, which linguistic variety is represented, how diagnostics and
documentation should sound, what compromises are acceptable for typing and
tooling, and how the project will avoid speaking for a community without
feedback from that community.

The project also shows why conventional semantics can be useful in this kind
of design study. If \plat{} also introduced an unusual type system, radical
control model, or experimental runtime, it would be harder to tell which design
effects came from the linguistic layer. By keeping the programming model
familiar, the artifact isolates regional-language identity as the main design
variable.

\begin{table}[t]
\centering
\caption{Design questions for culturally localized programming languages.}
\label{tab:checklist}
\begin{tabular}{p{0.32\linewidth}p{0.58\linewidth}}
\toprule
\textbf{Question} & \textbf{Why it matters} \\
\midrule
Who is the intended speech community? &
Clarifies whether the language is designed for fluent speakers, learners,
teachers, artists, researchers, or a broader public. \\
Which dialect or orthography is chosen, and why? &
Makes visible the linguistic authority and compromise behind keyword,
diagnostic, and documentation choices. \\
Which parts of the system are localized? &
Distinguishes surface syntax from diagnostics, documentation, examples,
tooling, libraries, editor support, and community materials. \\
How are authenticity and typability balanced? &
Identifies tradeoffs involving diacritics, ASCII compatibility, keyboard
layouts, searchability, and source portability. \\
Is the goal education, cultural visibility, practical use, art, or research? &
Prevents the project from being evaluated against the wrong success criteria. \\
How can community feedback be incorporated? &
Creates a path for revision, accountability, and speaker participation. \\
\bottomrule
\end{tabular}
\end{table}

This checklist is intended to be reusable beyond \plat. It can guide the
design of other regional-language or minority-language programming systems,
especially those that want to be taken seriously without claiming technical
novelty in traditional programming-language theory.

% TODO: Validate this checklist by applying it to additional localized
% programming-language projects.
